\section{本章小结}
\label{sec:opt-con-summary}

章首的预算分配例可以把全章串起来。最优点$\bm x^\star=(3/2,1/2)^\top$的梯度不为零，但负梯度落在可行域的法锥中；由于约束仿射、目标凸，乘子$\lambda_1^\star=1/2$与$\bm x^\star$一起满足KKT条件，这已证明全局最优。同一个乘子给出对偶下界：$q(1/2)=1/4=f(\bm x^\star)$，间隙为零，任何人都可以用这一对数据核验结果。它也回答了预算变化的问题：预算增加$0.1$，最优值至多下降$\lambda_1^\star\times0.1=0.05$，实际下降$0.0475$；在活动集不变的区间内，乘子就是最优值对预算的导数。求解方面，这个问题可以用投影（三角形上的投影可以显式计算）、惩罚或乘子法、或者把它当作二次规划用活动集法或内点法求解；它是凸问题，不需要松弛。

全章的结论可以按职责归纳。最优条件刻画候选：变分不等式只用可行域的几何，KKT条件把它写成梯度与乘子，从局部最优推出乘子存在需要约束资格（如LICQ），而凸问题中满足KKT条件就是全局最优。对偶认证目标值：弱对偶无条件地给出下界，Slater条件等正则条件保证零间隙和乘子存在，但原始最优点是否存在仍需另证；实际报告结果时，可行上界与可信下界之差是最可靠的证书。敏感性解释模型的改变：强对偶下，乘子给出最优值的全局单侧界，只有在值函数可微处才是导数；最优解的导数需要固定活动集、严格互补以及可逆的KKT矩阵，活动集变化或解不唯一时局部描述失效。算法产生候选：表~\ref{tab:opt-con-methods}比较了各类方法依赖的结构和每轮代价。

\begin{table*}[!htbp]
  \centering\small
  \begin{tabularx}{\linewidth}{@{}lXXX@{}}
    \toprule
    方法 & 适用结构 & 每轮主要代价 & 可行性与保证 \\
    \midrule
    投影梯度 & 光滑目标，投影容易的闭凸集 & 梯度加一次投影 & 每步可行；凸时$O(1/T)$，强凸时线性收敛 \\
    条件梯度 & 光滑目标，线性子问题容易的紧凸集 & 梯度加一次线性子问题 & 每步可行；凸时$O(1/t)$，解在面内部时振荡变慢 \\
    镜像下降 & 几何非欧氏的集合（如单纯形） & 梯度加一次散度子问题 & 保持在参考函数定义域内部 \\
    二次惩罚 & 一般光滑约束 & 一个无约束子问题 & 有限罚参数下不可行；参数增大导致病态 \\
    精确惩罚 & 乘子存在或误差界成立 & 一个非光滑子问题 & 罚参数超过阈值后最小点可行 \\
    增广Lagrange & 等式约束（不等式经松弛变量） & 一个光滑子问题加乘子更新 & 固定罚参数即可收敛（凸时） \\
    ADMM & 两块可分目标、线性耦合 & 两个可分子问题加乘子更新 & 残差、目标值与乘子收敛（凸时）；适合中等精度与分布式 \\
    SQP & 光滑非线性约束 & 一个二次规划 & 解附近二次或超线性收敛；需评价函数保证全局行为 \\
    内点法 & 线性规划与锥规划 & 一个KKT线性系统 & 迭代严格可行或残差逐步消除；间隙由障碍参数控制 \\
    \bottomrule
  \end{tabularx}
  \caption{本章约束优化方法的比较。表中“保证”只概括正文已说明的结论及其前提，具体条件见相应小节。}
  \label{tab:opt-con-methods}
\end{table*}

选择方法时可以按三个问题依次判断。第一，约束是什么类型：凸集上投影或线性子问题是否便宜，约束是否非线性，是否有矩阵约束。第二，子问题代价多少：一次投影、一次线性最小化、一个无约束问题、一个二次规划，还是一个KKT线性系统。第三，需要什么精度和证书：低到中等精度、需要分布式计算时，一阶方法与ADMM合适；需要高精度和原始—对偶证书时，内点法与SQP更合适。

最后要区分等价建模与松弛。把凸二次规划写成二阶锥规划、用Schur补把平方约束写成矩阵不等式，都是等价改写，最优值与最优解不变。半正定松弛则扩大了可行集合，它给出原最优值的单侧界，松弛解还必须恢复成原问题的候选并按原约束检验。在报告结果时，模型数据的敏感性、算法的计算误差、松弛的间隙与恢复的损失，应当分别说明来源和大小。
